\section{Kiến trúc tổng thể của hệ thống}

Phần này trình bày kiến trúc ở mức hệ thống. Các chi tiết về mô hình dữ liệu, Rule Evaluation Engine, API, cache và bảo mật được tách sang các mục chuyên sâu phía sau để tránh lặp lại nội dung.

\subsection{Các thành phần chính}

Hệ thống được tổ chức thành sáu thành phần: Admin Dashboard, Backend API, Rule Evaluation Engine, PostgreSQL/Cache, JavaScript SDK và Demo Application. Dashboard phục vụ quản trị cấu hình; Backend API là ranh giới truy cập duy nhất; Rule Evaluation Engine quyết định trạng thái flag tại runtime; database lưu cấu hình và audit log; cache tăng tốc truy xuất snapshot; SDK và Demo Application minh họa cách ứng dụng khách tích hợp Feature Flag.

\begin{figure}[H]
    \centering
    \includegraphics[width=0.95\textwidth]{\detokenize{Images/mermaid/Kiến trúc tổng thể của hệ thống.png}}
    \caption{Kiến trúc tổng thể của hệ thống Feature Flag Platform}
    \label{fig:system-overall-architecture}
\end{figure}


\subsubsection{Admin Dashboard}

Dashboard cung cấp giao diện quản trị cho project, Feature Flag, rule, flag group, audit log, history, RBAC demo và thống kê evaluation. Dashboard không tự quyết định kết quả evaluation mà gửi yêu cầu đến Backend API, nhờ đó logic nghiệp vụ và quyền truy cập vẫn được kiểm soát tập trung ở phía máy chủ.

\subsubsection{Backend API}

Backend API gom toàn bộ management endpoint và Evaluation API trong một service NestJS. Các endpoint quản trị xử lý validation, phân quyền, persistence và audit log; Evaluation API phục vụ ứng dụng tích hợp với response nhất quán gồm trạng thái bật/tắt và reason code.

\subsubsection{Rule Evaluation Engine}

Rule Evaluation Engine là lớp domain thuần, độc lập với controller. Engine nhận evaluation request, sử dụng snapshot cấu hình đã được chuẩn hóa và trả về quyết định xác định. Thứ tự ưu tiên chi tiết được trình bày trong mục thiết kế engine, vì vậy mục này chỉ xác định vị trí của engine trong kiến trúc tổng thể.

\subsubsection{Database và cache}

PostgreSQL là nguồn dữ liệu chính cho project, environment, flag, rule, group, audit log và thống kê aggregate. Cache chỉ được xem là tầng tăng tốc cho evaluation snapshot; dữ liệu trong cache có thể bị loại bỏ mà không làm mất tính đúng đắn vì database vẫn là nguồn sự thật.

\subsubsection{JavaScript SDK}

JavaScript SDK đóng gói cách gọi Evaluation API, chuẩn hóa timeout, lỗi mạng và fallback fail-closed cho ứng dụng web. SDK giúp Demo Application không phụ thuộc trực tiếp vào chi tiết endpoint và là cơ sở để mở rộng SDK cho các ngôn ngữ khác.

\subsubsection{Demo Application}

Demo Application mô phỏng một hệ thống thương mại điện tử nhỏ, sử dụng SDK để quyết định hiển thị tính năng theo kết quả evaluation. Vai trò của ứng dụng này là chứng minh giá trị thực tế của nền tảng thay vì chứa logic quản trị riêng.

\subsection{Control plane và data plane}

Kiến trúc được phân chia theo hai luồng trách nhiệm. \textit{Control plane} quản lý cấu hình và lịch sử thay đổi; \textit{data plane} xử lý evaluation ở runtime. Cách phân chia này phù hợp với đặc điểm của Feature Flag: thao tác quản trị ít hơn nhưng yêu cầu evaluation có thể xuất hiện liên tục từ ứng dụng tích hợp \cite{launchdarkly_feature_flags,unleash_feature_flag_best_practices}.

\begin{figure}[H]
    \centering
    \includegraphics[width=0.88\textwidth]{\detokenize{Images/mermaid/Control Plane - Data Plane.png}}
    \caption{Phân tách control plane và data plane trong hệ thống}
    \label{fig:control-plane-data-plane}
\end{figure}


\subsubsection{Vai trò của control plane}

Control plane bao gồm Dashboard, management endpoints, persistence và audit logging. Mọi mutation quan trọng đều đi qua validation, kiểm tra quyền, ghi database và tạo audit entry append-only.

\subsubsection{Vai trò của data plane}

Data plane bao gồm Evaluation API, Rule Evaluation Engine và snapshot cache. Luồng này được tối ưu cho quyết định nhanh, xác định và fail-closed khi thiếu cấu hình hoặc phát sinh lỗi.

\subsubsection{Luồng dữ liệu giữa hai thành phần}

Người quản trị thay đổi cấu hình qua Dashboard; Backend API lưu thay đổi, ghi audit log và invalidation snapshot liên quan. Ứng dụng khách gửi evaluation context qua SDK hoặc gọi trực tiếp Evaluation API; engine đọc snapshot hiện tại và trả kết quả cho ứng dụng.

\subsubsection{Lợi ích của việc tách biệt hai plane}

Sự tách biệt giúp giảm coupling giữa giao diện quản trị và runtime evaluation, cho phép tối ưu hai luồng theo tải khác nhau, đồng thời làm rõ ranh giới bảo mật: dữ liệu quản trị và lịch sử thay đổi không cần phơi bày cho ứng dụng khách.

\subsection{Luồng quản trị Feature Flag}

Luồng quản trị tập trung vào việc biến ý định phát hành thành cấu hình có kiểm soát. Các bước dưới đây chỉ mô tả trình tự kiến trúc; cấu trúc bảng dữ liệu và endpoint cụ thể được trình bày ở các mục sau.

\subsubsection{Tạo và chỉnh sửa Feature Flag}

Người quản trị khai báo flag key, tên và mô tả. Backend kiểm tra ràng buộc duy nhất trong project, tạo hoặc cập nhật bản ghi liên quan và trả kết quả cho Dashboard.

\subsubsection{Cấu hình trạng thái và serving mode}

Sau khi có flag, người quản trị cấu hình lifecycle/status và serving mode theo từng environment. Thiết kế này tách trạng thái quản lý của flag khỏi cách flag được phục vụ ở runtime.

\subsubsection{Cấu hình và sắp xếp rule}

Khi flag chạy ở chế độ targeted, người quản trị cấu hình các rule hỗ trợ bởi engine. Dashboard chỉ biểu diễn và gửi cấu hình; việc diễn giải rule thuộc trách nhiệm của Rule Evaluation Engine.

\subsubsection{Lưu audit log và invalidation cache}

Sau mỗi mutation, Backend API ghi audit log và làm mới snapshot liên quan. Nhờ vậy, hệ thống vừa có khả năng truy vết vừa hạn chế việc data plane sử dụng cấu hình cũ.

\subsection{Luồng đánh giá Feature Flag}

Luồng evaluation bắt đầu từ ứng dụng khách và kết thúc bằng một quyết định có reason code. Evaluation API không thay đổi cấu hình, không ghi audit log quản trị và không yêu cầu Dashboard tham gia trong runtime.

\begin{figure}[H]
    \centering
    \includegraphics[width=0.95\textwidth]{\detokenize{Images/mermaid/Luồng đánh giá Feature Flag.png}}
    \caption{Luồng đánh giá Feature Flag từ ứng dụng khách đến Rule Evaluation Engine}
    \label{fig:evaluation-flow}
\end{figure}


\subsection{Luồng tích hợp ứng dụng khách}

Ứng dụng khách có thể gọi trực tiếp Evaluation API hoặc sử dụng JavaScript SDK. Trong mini project, Demo Application dùng SDK để giảm mã lặp ở phía client và thể hiện cách tích hợp có thể tái sử dụng.

\subsubsection{Khởi tạo evaluation context}

Ứng dụng tạo context từ các định danh không nhạy cảm như targeting key, user ID mẫu, role và organization ID mẫu. Context chỉ là đầu vào cho engine, không phải bằng chứng phân quyền.

\subsubsection{Gọi trực tiếp Evaluation API}

Cách gọi trực tiếp phù hợp cho kiểm thử hoặc hệ thống muốn tự quản lý HTTP client. Client gửi project key, environment key, flag key và context, sau đó tự xử lý response.

\subsubsection{Gọi thông qua JavaScript SDK}

SDK bao bọc cùng contract nhưng bổ sung cấu hình base URL, timeout và fallback. Nội dung triển khai SDK cụ thể được trình bày trong chương kết quả.

\subsubsection{Áp dụng kết quả trong ứng dụng}

Ứng dụng chỉ sử dụng trường kết quả để bật hoặc ẩn hành vi giao diện tương ứng. Việc này giữ cho logic sản phẩm tách khỏi cách cấu hình flag được quản trị.

\subsubsection{Xử lý khi Evaluation API không khả dụng}

Khi không nhận được kết quả hợp lệ, client áp dụng mặc định an toàn là tắt tính năng. Quy tắc này tránh việc lỗi hạ tầng vô tình mở tính năng chưa sẵn sàng.

\subsection{Lựa chọn công nghệ}

Các công nghệ được chọn theo tiêu chí phù hợp phạm vi mini project, dễ kiểm thử và đủ gần với một hệ thống thực tế.

\subsubsection{Lựa chọn NestJS cho Backend API}

NestJS hỗ trợ cấu trúc module rõ ràng, dependency injection, DTO validation và test bằng Jest, phù hợp với backend REST cần tách controller, service và domain logic.

\subsubsection{Lựa chọn React cho giao diện web}

React phù hợp để xây dựng Dashboard và Demo Application dạng SPA, tái sử dụng component và phản hồi nhanh với trạng thái cấu hình thay đổi.

\subsubsection{Lựa chọn PostgreSQL và Prisma}

PostgreSQL cung cấp giao dịch và kiểu JSONB cho rule parameters/audit snapshot; Prisma giúp mô hình hóa schema, migration và truy cập dữ liệu có kiểu.

\subsubsection{Lựa chọn cache in-memory hoặc Redis}

In-memory cache đủ cho demo một instance; Redis là hướng mở rộng khi cần chia sẻ snapshot giữa nhiều backend instance. Thiết kế giữ cache là tùy chọn để không phụ thuộc vào hạ tầng ngoài.

\subsubsection{Lựa chọn Docker và Docker Compose}

Docker Compose gom backend, frontend, demo và database vào một quy trình khởi chạy thống nhất, thuận tiện cho chấm điểm và trình diễn.

\subsubsection{Các phương án thay thế và lý do không lựa chọn}

Các lựa chọn như microservice tách riêng evaluation service, message broker hoặc frontend framework khác chưa cần thiết trong phạm vi mini project. Hệ thống ưu tiên kiến trúc đơn giản, có ranh giới rõ và có thể mở rộng sau.
